\documentclass[sigconf]{acmart}

\usepackage{booktabs,multirow}
%\usepackage{amssymb}

%%
%% \BibTeX command to typeset BibTeX logo in the docs
\AtBeginDocument{%
  \providecommand\BibTeX{{%
    Bib\TeX}}}

%% Rights management information.  This information is sent to you
%% when you complete the rights form.  These commands have SAMPLE
%% values in them; it is your responsibility as an author to replace
%% the commands and values with those provided to you when you
%% complete the rights form.
\setcopyright{none}
\copyrightyear{2026}
\acmYear{2026}
\acmDOI{}
%% These commands are for a PROCEEDINGS abstract or paper.
%% DSAIT4055 Web Science and Engineering (2025/26 Q4)
\acmConference[DSAIT4055 '26]{DSAIT4055 Web Science and Engineering}{Q4,
  2026}{TU Delft}
%%
%%  Uncomment \acmBooktitle if the title of the proceedings is different
%%  from ``Proceedings of ...''!
%%
%%\acmBooktitle{Woodstock '18: ACM Symposium on Neural Gaze Detection,
%%  June 03--05, 2018, Woodstock, NY}
\acmISBN{}

%%
%% Submission ID.
%% Use this when submitting an article to a sponsored event. You'll
%% receive a unique submission ID from the organizers
%% of the event, and this ID should be used as the parameter to this command.
%%\acmSubmissionID{123-A56-BU3}

%%
%% For managing citations, it is recommended to use bibliography
%% files in BibTeX format.
%%
%% You can then either use BibTeX with the ACM-Reference-Format style,
%% or BibLaTeX with the acmnumeric or acmauthoryear sytles, that include
%% support for advanced citation of software artefact from the
%% biblatex-software package, also separately available on CTAN.
%%
%% Look at the sample-*-biblatex.tex files for templates showcasing
%% the biblatex styles.
%%

%%
%% The majority of ACM publications use numbered citations and
%% references.  The command \citestyle{authoryear} switches to the
%% "author year" style.
%%
%% If you are preparing content for an event
%% sponsored by ACM SIGGRAPH, you must use the "author year" style of
%% citations and references.
%% Uncommenting
%% the next command will enable that style.
%%\citestyle{acmauthoryear}


%%
%% end of the preamble, start of the body of the document source.
\begin{document}

%%
%% The "title" command has an optional parameter,
%% allowing the author to define a "short title" to be used in page headers.
\title{Framing Generative AI Governance in Online News: A Timeline Analysis from ChatGPT’s Launch to the EU AI Act (2022–2026)}
\subtitle{DSAIT4055 Web Science and Engineering (2025/26 Q4)}

%%
%% The "author" command and its associated commands are used to define
%% the authors and their affiliations.
%% Of note is the shared affiliation of the first two authors, and the
%% "authornote" and "authornotemark" commands
%% used to denote shared contribution to the research.
\author{Brewen Couaran}
\email{bcouaran@tudelft.nl}
\affiliation{%
  \institution{TU Delft (6533973)}
  \city{Delft}
  \country{The Netherlands}
}

\author{Yuvraj Singh Pathania}
\email{y.s.pathania@student.tudelft.nl}
\affiliation{%
  \institution{TU Delft (6495044)}
  \city{Delft}
  \country{The Netherlands}
}

\author{Arjun Rajesh Nair}
\email{A.RajeshNair@student.tudelft.nl}
\affiliation{%
  \institution{TU Delft (6327184)}
  \city{Delft}
  \country{The Netherlands}
}

%%
%% By default, the full list of authors will be used in the page
%% headers. Often, this list is too long, and will overlap
%% other information printed in the page headers. This command allows
%% the author to define a more concise list
%% of authors' names for this purpose.
\renewcommand{\shortauthors}{Couaran et al.}

\received{\today}

%%
%% The code below is generated by the tool at http://dl.acm.org/ccs.cfm.
%% Please copy and paste the code instead of the example below.
%%
\begin{CCSXML}
<ccs2012>
    <concept>
        <concept_id>10010405.10010444.10010447</concept_id>
        <concept_desc>Applied computing~Computational social science</concept_desc>
        <concept_significance>500</concept_significance>
    </concept>
    <concept>
        <concept_id>10003451.10003452.10003453</concept_id>
        <concept_desc>Social and professional topics~Computing / technology policy</concept_desc>
        <concept_significance>300</concept_significance>
    </concept>
    <concept>
        <concept_id>10010405.10010407.10010408.10010409</concept_id>
        <concept_desc>Applied computing~World Wide Web</concept_desc>
        <concept_significance>300</concept_significance>
    </concept>
</ccs2012>
\end{CCSXML}

\ccsdesc[500]{Applied computing~Computational social science}
\ccsdesc[300]{Social and professional topics~Computing / technology policy}
\ccsdesc[300]{Applied computing~World Wide Web}

%%
%% Keywords. The author(s) should pick words that accurately describe
%% the work being presented. Separate the keywords with commas.
\keywords{Generative AI, AI Governance, Online News Framing, Computational Social Science, Web Science, ChatGPT, GDELT 2.0}

\begin{abstract}
The public release of ChatGPT in late 2022 rapidly turned generative AI from a technological novelty into a central governance concern, prompting legislation, international summitry, and sustained public debate about regulation, safety, and accountability~\cite{cheung_generative_2025,taeihagh_governance_2025}.
Online news plays a central role in shaping how these issues are publicly understood, making web-based news infrastructure a critical object of Web Science inquiry~\cite{lazer_computational_2009}.

This paper presents one of the first large-scale longitudinal framing analyses of generative AI governance in online news.
Drawing on 1,116,091 articles from the GDELT~2.0 Global Knowledge Graph~\cite{kalev_gdelt_2013}, spanning November~2022 to June~2026, we apply a six-category governance frame taxonomy (Innovation \& Opportunity, Risk \& Safety, Regulation \& Governance, Rights \& Privacy, Economic Competition \& Labour, and Misinformation \& Integrity) via a two-stage procedure combining multilingual keyword matching and LaBSE sentence-embedding confirmation~\cite{kwak_frameaxis_2021,feng_language-agnostic_2022}.

We document a discourse arc from optimism to precaution: Innovation framing peaks in February~2023 before retreating, while Risk \& Safety rises steadily from 2024 onward.
An event study across eight governance milestones shows that media framing responds selectively and unevenly to policy events: high-profile intergovernmental summits coincide with less risk framing but more governance vocabulary, civil-society alarm signals raise risk framing without any governance-response coverage, and routine administrative deadlines produce no lasting shift.
Regional analysis shows US outlets are at once the most regulation-intensive and innovation-positive, while EU outlets carry proportionally higher risk framing relative to their regulation share~\cite{nguyen_generative_2026}.
These findings contribute to Web Science and computational social science by showing how web-based news infrastructures shape public discourse around the governance of emerging AI technologies.
\end{abstract}

%%
%% This command processes the author and affiliation and title
%% information and builds the first part of the formatted document.
\maketitle

\section{Introduction}

Between November~2022 and mid-2026, online news coverage of generative AI changed character. The early post-ChatGPT optimism, focused on breakthroughs and market potential, gave way to a more anxious debate about risk, regulation, and global competition~\cite{taeihagh_governance_2025,cheung_generative_2025}.
Within months of ChatGPT's public release, policymakers on several continents were holding summits and drafting legislation, most visibly the EU AI Act~\cite{gstrein_general-purpose_2024}.
News coverage was central to this shift. It is the main channel through which most citizens encounter policy debates, so the online news ecosystem both mirrored and helped shape public understanding of what AI governance means and why it matters~\cite{lazer_computational_2009}.

How that coverage is framed, which concerns get emphasized and who is held responsible, is therefore a Web Science question with direct policy relevance.
Media framing theory holds that the aspects of a topic journalists choose to highlight shape how audiences interpret it~\cite{entman_framing_1993}. Prior computational work has found AI coverage to be persistently more optimistic than pessimistic, even as concerns about job displacement and loss of control have grown over time~\cite{fast_long-term_2017}.
Yet most existing analyses of AI in the news rely on small, single-country samples or short observation windows~\cite{chuan_framing_2019,eisenbeis_ai_2024}.
More recent studies of \emph{generative} AI coverage~\cite{hendrickx_tool_2026,nguyen_generative_2026,lee_framing_2026} do address the post-ChatGPT era, but they focus on broad AI narratives or on elite outlets in a single country. The global governance discourse, spanning the full policy life-cycle from ChatGPT's launch to the EU AI Act's GPAI obligations, has not been mapped at this scale.

This paper addresses that gap.
We present a large-scale computational analysis of 1,116,091 online news articles drawn from the GDELT~2.0 Global Knowledge Graph~\cite{kalev_gdelt_2013} via Google BigQuery, spanning November~2022 to June~2026.
Applying a six-category governance frame taxonomy informed by established framing typologies~\cite{semetko_framing_2000}, we measure how the volume and frame composition of generative AI governance evolved month by month, and examine whether eight major AI governance milestones, from the ChatGPT launch to the EU AI Act's GPAI obligations, coincide with discernible shifts in coverage or framing.

We address three research questions:
\begin{enumerate}
    \item[\textbf{RQ1.}] How has the monthly volume of online news coverage of generative AI governance changed between late 2022 and mid-2026?
    \item[\textbf{RQ2.}] Which governance-related frames dominate coverage, and how does their relative salience evolve over time?
    \item[\textbf{RQ3.}] Do major AI governance milestones coincide with noticeable shifts in coverage volume or dominant framing?
\end{enumerate}

We make three contributions.
We provide one of the largest longitudinal framing analyses of generative AI governance in online news to date.
We document a shift from innovation- and market-oriented framing (2022--2023) toward risk- and regulation-dominated framing (2024--2026), with a geopolitically charged acceleration in 2025.
And we show that media responsiveness to governance events is selective and uneven: high-profile intergovernmental summits coincide with less risk framing but more governance vocabulary, civil-society alarm signals raise risk framing without any matching governance-response coverage, and routine administrative deadlines produce no lasting shift at all. News agenda-setting, in other words, seems to track political contestation more than the legal calendar.

\section{Background and Related Work}

\subsection{Media Framing Theory}

Framing theory holds that communicators select and emphasize certain aspects of a perceived reality, thereby promoting particular problem definitions, causal interpretations, and remedies~\cite{entman_framing_1993}.
Semetko and Valkenburg~\cite{semetko_framing_2000} operationalized this by deriving five generic frames (attribution of responsibility, conflict, economic consequences, human interest, and morality) from European political news, establishing the standard typology that AI-framing studies have since inherited.
Our governance-specific taxonomy follows the same logic: each category maps a distinct evaluative lens through which journalists interpret AI regulation.
DiMaggio et al.~\cite{dimaggio_exploiting_2013} further demonstrated that dictionary-based and probabilistic methods can reliably recover latent frames in large newspaper corpora, directly validating the scalability of our approach.

\subsection{Computational Analysis of AI in the News}

Lazer et al.~\cite{lazer_computational_2009} established computational social science as a paradigm for extracting insight from digital trace data at scales infeasible with manual coding.
The GDELT~2.0 Global Knowledge Graph~\cite{kalev_gdelt_2013}, updated every fifteen minutes across over one hundred languages, provides structured metadata (topics, themes, tone, and geolocation) for news articles worldwide, making it well-suited for large-scale extraction via BigQuery.

Fast and Horvitz~\cite{fast_long-term_2017} applied this paradigm to thirty years of \emph{New York Times} AI coverage, finding that volume grew sharply after 2009 and tone remained consistently optimistic even as concerns about job displacement and loss of control intensified.
Chuan et al.~\cite{chuan_framing_2019} subsequently showed that business and economic frames dominate AI news, and that risks, though discussed less frequently than opportunities, are articulated with greater specificity.
These findings establish the pre-ChatGPT baseline against which our 2022--2026 corpus can be contextualized.

\subsection{Generative AI Coverage and Governance}

The release of ChatGPT prompted a new wave of empirical attention to AI in the media.
Eisenbeis and Elkins~\cite{eisenbeis_ai_2024} found US AI news sentiment remained predominantly positive in its immediate aftermath, though governance-specific coverage was underexplored.
Hendrickx and Van Coppenolle~\cite{hendrickx_tool_2026} offer the methodologically closest comparison: analyzing a Belgian quality newspaper (2020--2024), they identify twelve frames distributed roughly equally between tool and threat narratives across 516~articles.
Nguyen and Wischnewski~\cite{nguyen_generative_2026} examined 18,032 articles from fifteen international outlets, finding that anthropomorphic framing characterizes fewer than nine percent of texts, that economic themes dominate, and that European outlets are more governance-oriented than their North American or Asian counterparts.
Lee and Park~\cite{lee_framing_2026} analyzed the \emph{New York Times} and \emph{Wall Street Journal}, identifying societal and economic impact as the most prevalent frames and responsibility attributed predominantly at the societal rather than individual level.

On the governance side, Cheung~\cite{cheung_generative_2025} traces how crisis framing triggered a strategic pivot in UK AI policy; Taeihagh~\cite{taeihagh_governance_2025} enumerates six governance challenges, from data governance to power imbalances, that map directly onto our frame taxonomy; and Gstrein et al.~\cite{gstrein_general-purpose_2024} characterizes the EU AI Act as a proactive governance benchmark with global reach.
These studies establish useful precedents, but each is bounded by national scope, a short observation window, a broad AI focus rather than governance-specificity, or the absence of a policy-milestone design.
Our study combines global GDELT coverage, a 44-month governance-specific corpus, and a structured event study, extending prior work on all three fronts at once.

\section{Dataset and Methodology}

\subsection{Corpus Construction}

We build our corpus from the GDELT~2.0 Global Knowledge Graph (GKG)~\cite{kalev_gdelt_2013}, accessed via Google BigQuery.
The GKG indexes online news articles at global scale, providing structured metadata per record: named entities (\texttt{AllNames}), quotation text (\texttt{Quotations}), thematic codes (\texttt{V2Themes}), automated tone (\texttt{V2Tone}), and geolocation (\texttt{V2Locations}).

An article enters our corpus when it meets two conditions at once.
First, it must carry a generative AI signal: a 34-term lexicon (including \textit{ChatGPT}, \textit{large language model}, \textit{GPT-4}, \textit{Gemini}, \textit{Anthropic}, \textit{DeepSeek}, \textit{AI agent}) matched against \texttt{AllNames}, with a multilingual contextual subset additionally matched against \texttt{Quotations}.
Second, it must carry a governance signal, operationalized as any of three complementary indicators: 
\begin{enumerate}
    \item 27 regulation and policy keywords in \texttt{Quotations}.
    \item 17 GDELT thematic codes spanning legal, accountability, and policy categories (e.g. \texttt{ECON\_REGULATION}, \texttt{LEGISLATION}, \texttt{WB\_831\_GOVERNANCE}).
    \item 18 URL-slug patterns (e.g. \texttt{-regulat}, \texttt{-policy-}, \texttt{-ai-act}).
\end{enumerate}

No language or source filter is applied, preserving GDELT's global reach.
Querying from 2022-11-01 to 2026-06-30 yields 1,116,091 articles across 44~months after deduplication.
As a sanity check, governance-flagged articles make up 50.7\% of all GenAI-matching articles over the same period, so the filter captures a large, reasonably well-bounded slice of GenAI discourse.

\subsection{Frame Taxonomy and Assignment}

We operationalize governance framing through a six-category taxonomy (Table~\ref{tab:frames}), designed to capture the principal evaluative lenses through which journalists interpret AI regulation.
The categories are informed by Semetko and Valkenburg's~\cite{semetko_framing_2000} generic frame typology and adapted to the generative AI domain following recent media studies.
Each frame is defined by a multilingual keyword dictionary across nine languages (English, French, German, Spanish, Italian, Dutch, Polish, Portuguese, and Romanian): 30--68 English terms per frame (318 in total), expanded to 959 terms across all languages.
To guard against substring false-positives, single-word terms are wrapped in word-boundary anchors; multi-word phrases are left as plain substrings given their inherent specificity.

Frame assignment uses a two-stage procedure.
In the first stage, keyword matching counts hits for each frame across the concatenated \texttt{AllNames} and \texttt{Quotations} fields, approximating named-entity and salient-phrase coverage.
In the second stage, we score every article with a bipolar sentence-embedding model (\texttt{LaBSE}, Language-Agnostic BERT Sentence Embeddings~\cite{feng_language-agnostic_2022}, covering 109 languages): for each frame, we compute the cosine similarity between the article embedding and the frame's positive-pole keyword centroid minus the negative-pole centroid~\cite{kwak_frameaxis_2021}.
A frame is \emph{confirmed} for an article only when both the keyword flag fires \emph{and} the embedding score is positive, i.e.\ the article text is semantically closer to the positive pole than to the negative pole of the frame axis~\cite{widmann_creating_2023}.
This intersection filters incidental keyword matches while remaining interpretable: the confirmed label requires both formal vocabulary and semantic alignment.

The scheme is multi-label: an article can be confirmed in several frames at once.
Overall, 40.8\% of articles (455,349) are confirmed in at least one frame; the rest contribute to volume counts but are excluded from frame-share analyses.
The dominant frame for per-article analyses is the one with the highest keyword hit count (normalized by dictionary size) among confirmed frames.

\begin{table}[!t]
\caption{Governance frame taxonomy with example keywords.}
\label{tab:frames}
\footnotesize
\setlength{\tabcolsep}{4pt}
\begin{tabular}{p{1.9cm}p{2.5cm}p{2.4cm}}
\toprule
\textbf{Frame} & \textbf{Description} & \textbf{Example keywords} \\
\midrule
Innovation \& Opportunity & Benefits, progress, transformative potential & innovation, breakthrough, opportunity, efficiency \\[2pt]
Risk \& Safety & Harms, threats, security risks, existential concerns & harm, threat, risk, misuse, hallucination \\[2pt]
Regulation \& Governance & Laws, oversight, compliance, institutional steering & regulation, policy, law, compliance, AI Act \\[2pt]
Rights \& Privacy & Data protection, civil liberties, copyright, fairness & privacy, copyright, bias, discrimination \\[2pt]
Econ.\ Comp.\ \& Labour & Job impacts, market dynamics, corporate competition & jobs, automation, market, investment \\[2pt]
Misinfo.\ \& Integrity & Deepfakes, disinformation, election integrity & deepfake, disinformation, election \\
\bottomrule
\end{tabular}
\end{table}

A methodological note: the Regulation \& Governance frame dictionary shares 14 of its English terms (and their cross-language equivalents, 89 of 183 terms in all) with the corpus-level governance filter.
As a robustness check we recomputed the frame with these shared terms removed: prevalence falls from 16.6\% to 13.9\%, but Regulation remains the most prevalent frame (ahead of Risk \& Safety at 12.2\%), and every milestone delta changes by less than 1~pp.
The embedding confirmation step further reduces the structural inflation by rejecting articles whose text does not align semantically with the frame.
Cross-frame comparisons should still be read as relative emphasis rather than absolute prevalence, but the frame's dominance is not an artefact of the shared vocabulary alone.

\subsection{Tone Measurement}

GDELT computes an automated tone score (\texttt{V2Tone}) for each GKG record based on lexical cues in the article text.
We extract the first component of this field, where positive values indicate more favorable coverage and negative values more critical coverage.
Tone is aggregated monthly across all corpus articles and separately by dominant frame to examine affective variation across governance narratives.
Given that GDELT's tone scoring is rule-based and susceptible to noise, we treat tone results as supplementary and indicative throughout.

\subsection{Event Study Design}

To address RQ3, we adopt a before--after event-study design.
For each of eight major AI governance milestones, we define an adaptive symmetric window: for each milestone the window length equals $\min(d_{\mathrm{prev}},\, d_{\mathrm{next}},\, 90)$~days, where $d_{\mathrm{prev}}$ and $d_{\mathrm{next}}$ are the distances to the preceding and following milestones respectively.
This prevents adjacent milestones from contaminating each other's comparison periods (e.g.\ the Bletchley Summit and the EU AI Act agreement, only 37~days apart, each receive a 37-day symmetric window).
We measure the change in regulation and risk frame prevalence between the pre-milestone and post-milestone periods.
We test each change two ways: a label-permutation test (10{,}000 random pre/post reassignments within the window) checks whether a delta exceeds sampling noise, and a window-matched placebo, the same delta computed at random non-milestone dates, gauges how large it is relative to ordinary temporal variation.
The milestones span the full governance lifecycle: the ChatGPT launch (November 2022), the Pause AI open letter (March 2023), the Bletchley Summit (November 2023), the EU AI Act political agreement (December 2023), the Seoul AI Safety Summit (May 2024), the EU AI Act's entry into force (August 2024), the OpenAI o1 public release (September 2024), and the GPAI model obligations (August 2025).
Three additional recent events (DeepSeek-R1 and the Trump AI executive order in January 2025, EU AI Act prohibited-practices enforcement in February 2025, and the Claude Fable~5 release with US AI export controls in June 2026) yield comparison windows of 8--13~days and are noted in the discussion as indicative rather than conclusive.

\subsection{Limitations}

A few limitations should be kept in mind.
Although the frame dictionaries cover nine languages, GDELT's indexing over-represents English-language sources; the multilingual terms extend coverage in principle, but in practice the corpus skews Anglophone, limiting generalizability to non-Anglophone media ecosystems.
This also confounds the regional comparison (§4.4): because our EU subset is restricted to English-language sources with EU-registered domains, the US--EU framing gap may partly reflect GDELT's language and sourcing composition rather than genuine editorial differences.
Frame assignment relies on \texttt{AllNames} and \texttt{Quotations} rather than full article prose: \texttt{AllNames} fires on named entities throughout the article and \texttt{Quotations} captures direct source speech from anywhere in the body, so coverage is article-wide rather than headline-only, but continuous journalistic text between named entities and quotes is not analysed; subtler framing that is neither entity-driven nor quote-driven may be missed.
The vocabulary overlap between the corpus governance filter and the Regulation frame inflates that frame's base rate, although the robustness check in §3.2 leaves its rank and the milestone deltas essentially unchanged.
The event-study deltas, while mostly significant against sampling noise, are modest relative to background temporal volatility (§4.4); milestone attributions are therefore suggestive rather than causal.
The temporal proximity of the Bletchley Summit and the EU AI Act agreement complicates the attribution of framing shifts to either event in isolation.
Finally, GDELT's \texttt{V2Tone} metric is rule-based and may not capture rhetorical nuance; tone results are treated as indicative throughout.

\subsection{Data and Code Availability}
All BigQuery extraction queries, multilingual frame dictionaries, Python frame-assignment pipeline, and analysis notebooks are publicly available at \url{https://github.com/brewcoua/GenAI-GDELT}.
The article-level framing dataset (1,116,091 articles with binary keyword flags and LaBSE embedding scores for all six frames) is published on HuggingFace at \url{https://huggingface.co/datasets/brewcoua/genai-gdelt-framing}.
A companion website with figures and citation is available at \url{https://brewcoua.github.io/GenAI-GDELT}.

\section{Analysis and Results}

\subsection{Coverage Overview}

Monthly volume grew from 14,908--17,788 articles in late 2022 to 20,680--28,582 articles by spring 2026, an approximately 59\% increase (1.59$\times$) over the study window.
The governance filter is inclusive by design: 50.7\% of all GDELT-indexed GenAI articles carry a governance signal, leaving 49.3\% as general-interest GenAI coverage excluded from the corpus.
Of the 1,116,091 retained articles, 40.8\% (455,349) are confirmed in at least one frame; the remaining 59.2\% contribute to volume counts but are not assigned a frame.

\subsection{RQ1: Monthly Volume Trends}

\begin{figure}[t]
  \includegraphics[width=\columnwidth]{figures/fig_1_volume_frames}
  \caption{Daily confirmed frame prevalence (14-day rolling mean, stacked area, left axis) and total articles per day (right axis), with governance-milestone markers.
           Regulation \& Governance dominates throughout; Risk \& Safety rises sharply from 2024 onward as Innovation \& Opportunity retreats after its February~2023 peak; the elevated plateau from January~2025 reflects the DeepSeek-R1 and Trump AI executive-order period.}
  \label{fig:vol_frames}
\end{figure}

Figure~\ref{fig:vol_frames} shows that coverage grew steadily from the ChatGPT launch onward, with no single month of dramatic collapse.
The growth trajectory plateaus at 22,000--26,000 articles per month through 2023--2024 before resuming a modest upward trend in 2025--2026.
Volume alone does not discriminate between sustained journalistic interest and shifting governance salience; the frame decomposition addresses this.

\subsection{RQ2: Frame Distribution and Temporal Evolution}

Table~\ref{tab:frame_stats} summarizes confirmed frame prevalence across the full corpus.
Regulation \& Governance is the most prevalent frame, at 16.6\% of all articles (185,245 confirmed), followed by Risk \& Safety (12.2\%, 136,379) and Innovation \& Opportunity (9.6\%, 106,708).
Economic Competition \& Labour (7.7\%, 85,657) is a close fourth. Rights \& Privacy and Misinformation \& Integrity stay marginal at 3.5--3.6\% each, a niche presence in governance coverage.

\begin{table}[!t]
\caption{Confirmed frame prevalence (keyword $\cap$ embedding). \% Corpus: fraction of all 1,116,091 articles confirmed in the frame. Monthly (mean): mean monthly prevalence rate across 44 months.}
\label{tab:frame_stats}
\footnotesize
\setlength{\tabcolsep}{3pt}
\begin{tabular}{lrrr}
\toprule
\textbf{Frame} & \textbf{Articles} & \textbf{\%~Corpus} & \textbf{Monthly} \\
\midrule
Regulation \& Gov.     & 185,245 & 16.6 & 16.5\% \\
Risk \& Safety         & 136,379 & 12.2 & 12.0\% \\
Innovation \& Opp.     & 106,708 &  9.6 &  9.6\% \\
Econ.\ Comp.\ \& Lab.  &  85,657 &  7.7 &  7.7\% \\
Rights \& Privacy      &  40,055 &  3.6 &  3.6\% \\
Misinfo.\ \& Integrity &  39,236 &  3.5 &  3.5\% \\
\bottomrule
\end{tabular}
\end{table}

The temporal dimension of the confirmed frame series (Figure~\ref{fig:vol_frames}) reveals a discourse shift.
From late 2022 to mid-2023, Regulation \& Governance dominates; Innovation \& Opportunity peaks in February~2023 (max 13.1\%), reflecting the initial wave of post-ChatGPT enthusiasm, with a secondary local peak (\textasciitilde11--12\%) around the Bletchley Summit in November~2023.
From 2024 onward, Risk \& Safety framing rises steadily, reaching its corpus-period maximum of 20.3\% in isolated months, while Innovation \& Opportunity settles below 10\% and remains there.
This reorientation from opportunity to risk framing mirrors the real-world trajectory of frontier AI debate: the post-ChatGPT enthusiasm giving way to expert concern about model capabilities and safety.

As an independent check, GDELT's automated tone score tracks the frame labels: Innovation articles are the only positively-valenced group (mean V2Tone $+1.11$), while every governance-adjacent frame is negative and Risk \& Safety is the most critical ($-2.30$). Read with the usual caution about a rule-based metric, this suggests governance coverage carries a critical register~\cite{fast_long-term_2017}.

\subsection{RQ3: Governance Milestones and Framing Shifts}

Table~\ref{tab:eventstudy} reports the change in confirmed Regulation and Risk frame prevalence across the adaptive symmetric window for each milestone.
For the ChatGPT launch, the 29-day window uses the first 29 days of corpus data as the pre-period (November 1--29, 2022), before the public release on November 30.

\begin{table}[!t]
\caption{Event-study results: change in confirmed frame prevalence (adaptive symmetric window, in percentage points; see §3.4). Positive = frame more prevalent post-milestone. $^{*}p<0.05$, $^{**}p<0.01$ from a 10{,}000-iteration label-permutation test within each window; unmarked deltas are not distinguishable from chance reassignment. Note that significance reflects the large per-window sample: most deltas remain within the range of ordinary temporal variation (see §4.4).}
\label{tab:eventstudy}
\footnotesize
\setlength{\tabcolsep}{3pt}
\begin{tabular}{lcrr}
\toprule
\textbf{Milestone} & \textbf{Window} & \textbf{$\Delta$ Reg.} & \textbf{$\Delta$ Risk} \\
\midrule
ChatGPT launch (Nov 2022)        & 29d  & $-1.7^{**}$ & $+0.6^{*}$ \\
Pause AI letter (Mar 2023)       & 90d  & $+0.4$ & $+1.5^{**}$ \\
Bletchley Summit (Nov 2023)      & 37d  & $+1.1^{**}$ & $-4.9^{**}$ \\
EU AI Act agreement (Dec 2023)   & 37d  & $+2.5^{**}$ & $+3.1^{**}$ \\
Seoul Summit (May 2024)          & 72d  & $-0.2$ & $-0.8^{**}$ \\
EU AI Act in force (Aug 2024)    & 42d  & $+1.5^{**}$ & $+3.0^{**}$ \\
OpenAI o1 release (Sep 2024)     & 42d  & $-1.4^{**}$ & $-0.1$ \\
GPAI obligations (Aug 2025)      & 90d  & $-1.9^{**}$ & $-0.3$ \\
\bottomrule
\end{tabular}
\end{table}

The results reveal a selective and asymmetric pattern, though the magnitudes are modest.
The permutation test confirms that most deltas are not sampling artefacts (Table~\ref{tab:eventstudy}), but because each window holds thousands of articles, even sub-point shifts clear significance; the more demanding comparison is against background volatility.
At random non-milestone dates, same-length windows already swing by several percentage points at the 95\% level (up to roughly $\pm5$--$8$~pp for the shorter windows), and every individual milestone delta falls within this band of ordinary period-to-period variation.
We therefore read the event study as evidence of a relative pattern across milestones rather than proof that any single milestone reshaped coverage.
The Bletchley Summit ($+1.1$~pp regulation) coincides with a sharp \emph{drop} in risk framing ($-4.9$~pp). One reading is that the summit's coverage recast existential-risk talk as governance-response language, with policy vocabulary crowding out alarm.
The EU AI Act political agreement, by contrast, lifted \emph{both} frames together ($+2.5$~pp regulation, $+3.1$~pp risk), matching coverage that framed legislative finalization as both a governance achievement and a reminder of continuing hazard.
The Pause AI open letter coincides with a Risk rise ($+1.5$~pp) and almost no Regulation change ($+0.4$~pp), which fits the idea that civil-society alarm signals raise risk framing without prompting governance-solution coverage.
The Seoul Summit ($-0.2$~pp, $-0.8$~pp) and OpenAI o1 release ($-1.4$~pp, $-0.1$~pp) produce near-zero deltas (the Seoul regulation and o1 risk changes are not statistically significant), pointing to lower political salience: the Seoul Summit was widely seen as reaffirming rather than advancing the Bletchley agenda, and a capability-oriented product release does not carry the legislative framing cues that drive governance discourse.
The EU AI Act's entry into force yields modest positive deltas ($+1.5$~pp regulation, $+3.0$~pp risk), and GPAI obligations produce a small negative regulation delta ($-1.9$~pp) with no significant risk movement ($-0.3$~pp), suggesting that by August 2025 the EU policy frame had become routine background.
Three additional events have comparison windows too short for reliable inference: DeepSeek-R1 combined with the Trump AI executive order (13~days), EU Article~5 prohibited-practices enforcement (13~days), and the Claude Fable~5 release with US export controls (8~days). All three coincide with large positive regulation spikes visible in Figure~\ref{fig:vol_frames}, but the window lengths preclude causal attribution; they are noted as indicative of dominant direction only.
Overall, only politically contested, novel, or internationally coordinated milestones generate discernible framing shifts; routine administrative deadlines do not.

\subsubsection{Regional variation}
Confirmed frame prevalence differs meaningfully across geographies (Figure~\ref{fig:regional}).
US outlets show the highest regulation prevalence (22.9\%) and innovation prevalence (10.1\%); this fits a period of contentious Congressional AI debates and the Biden executive order, both of which generated dense regulatory vocabulary.
EU outlets (English-language sources with EU-registered domains in GDELT; EU-language media is likely underrepresented) show relatively higher risk/safety prevalence (8.8\%) for their regulation share (13.5\%), which echoes Nguyen and Wischnewski's~\cite{nguyen_generative_2026} finding that European outlets are more governance-oriented in tone.
UK outlets sit between the two (16.8\% regulation, 13.9\% risk), reflecting both proximity to EU policy and high coverage of the Bletchley Summit.

The temporal evolution of these regional differences is visible in Figure~\ref{fig:regional}.
US regulation framing diverges sharply upward from 2024 onward, tracking the Biden executive order on AI, the OpenAI o1 period, and the legislative flurry of 2025; EU regulation framing, while lower in absolute terms, shows a parallel rise that plateaus later.
Risk framing in the UK tracks the US closely through 2024 before converging toward EU levels in 2025, matching the UK's post-Bletchley pivot from summitry to domestic AI-governance legislation.

\begin{figure}[t]
  \includegraphics[width=\columnwidth]{figures/fig_2_regional_framing}
  \caption{Daily confirmed frame prevalence for US, EU, and UK outlets (14-day rolling mean) across three governance frames.
           US regulation framing is consistently the highest and diverges sharply upward from 2024~onward; EU risk framing tracks closely with the UK throughout the study period.
           Vertical dashed lines mark the eleven governance milestones.}
  \label{fig:regional}
\end{figure}

\section{Discussion}

Our results show how generative AI governance discourse shifted alongside the real-world trajectory of AI policy.
The post-ChatGPT window (late 2022 to mid-2023) was dominated by regulatory vocabulary as legislators worldwide drafted responses to a novel and rapidly diffusing technology, with an innovation-framing peak in February~2023 reflecting the initial wave of enthusiasm.
The 2023 Bletchley period brought a secondary innovation uptick alongside sustained regulation coverage, reflecting the simultaneous excitement around frontier models and the urgency of international coordination.
From 2024 onward, Risk \& Safety framing grows relative to innovation, and the discourse shifts to a more sober, geopolitically charged tone, accelerating through 2025 with the DeepSeek-R1 and Trump executive-order period.
This pattern aligns with Taeihagh's~\cite{taeihagh_governance_2025} taxonomy of governance challenges: concerns about accountability, safety, and power asymmetries become more salient as capabilities advance past the initial novelty phase.

The event-study results offer a qualification to the volume-based narrative, with the caveat that the shifts are modest relative to background variation (§4.4).
Read as a relative pattern across milestones, framing appears to track political contestation and symbolic mobilization more than the legal calendar.
The Bletchley Summit nudges regulation framing upward ($+1.1$~pp) while risk framing drops ($-4.9$~pp), as if the event's diplomatic framing, centered on governance solutions and coordination, pushed alarm-oriented coverage aside.
This contrasts with the EU AI Act political agreement, where both regulation and risk framing rose together ($+2.5$ and $+3.1$~pp respectively): legislative finalization generated both policy-process coverage and renewed attention to the hazards the law was designed to address.
The Pause AI open letter, by contrast, amplified risk without a corresponding regulation response ($+1.5$ and $+0.4$~pp), which suggests that civil-society alarm signals and government governance signals are handled differently by news outlets.
The EU AI Act's administrative deadlines, specifically entry into force ($+1.5$~pp regulation) and GPAI obligations ($-1.9$~pp), produce modest or negative shifts despite being the most legally consequential events in the dataset.
This pattern is in line with intermittent agenda-setting theory~\cite{lazer_computational_2009}: journalists prioritize conflict, novelty, and elite political mobilization over procedural milestones, even when those milestones carry the weight of binding law.

The 2025 events extend this pattern to a new phase of the governance lifecycle.
Although the comparison windows for the DeepSeek-R1 / Trump AI executive order period and the EU Article~5 prohibited-practices enforcement are too short for robust inference, both show large positive regulation spikes visible in Figure~\ref{fig:vol_frames} as the elevated Regulation \& Governance plateau from January 2025 onward.
This suggests that the post-2024 governance discourse became more contentious, more geopolitically charged, and more densely regulation-framed, driven by the convergence of frontier capability announcements, US executive action, and EU enforcement milestones.

The regional picture adds another piece.
US outlets are at once the most regulation-intensive and the most innovation-positive, which fits a media ecosystem covering both a contentious bipartisan AI policy debate and the world's largest AI industry.
EU outlets combine lower innovation framing with proportionally higher risk framing for their level of regulation coverage, echoing prior evidence that European media frames AI more cautiously~\cite{nguyen_generative_2026} and reflecting the EU AI Act's foundational premise that AI poses fundamental rights risks requiring precautionary governance~\cite{gstrein_general-purpose_2024}.

Two methodological caveats deserve repeating.
First, even after embedding confirmation, the Regulation \& Governance frame stays inflated by the vocabulary it shares with the corpus filter; cross-frame comparisons measure relative emphasis, not absolute salience.
Second, our two-stage frame assignment trades recall for precision. The 59.2\% of unframed articles are mostly coverage that our closed vocabulary does not match, lexically or semantically, rather than noise, which is a limitation of any closed-vocabulary approach.
A human-annotated gold standard (e.g., via LLM-based labeling as in Gilardi et al.~\cite{gilardi_chatgpt_2023} and Ziems et al.~\cite{ziems_can_2024}) would allow quantitative error estimation and remains the most credible next validation step.

\section{Conclusion}

We presented one of the first large-scale, longitudinal framing analyses of generative AI governance news, covering 1,116,091 articles over 44 months and six frames with a two-stage keyword and multilingual sentence-embedding procedure.
Across the study period, Regulation \& Governance framing stays dominant, but its surroundings change: Innovation \& Opportunity peaks in February~2023 and then retreats, while Risk \& Safety climbs from 2024 onward, moving the discourse from early optimism toward precaution in step with the frontier AI debate.
The milestone response, read as a relative pattern rather than causal effect, is uneven and selective: risk framing fell around the Bletchley Summit as governance vocabulary rose, the EU AI Act agreement lifted both frames together, and the Pause AI letter raised alarm with no matching governance-response signal.
The events that moved coverage were the contested or novel ones; the routine administrative deadlines that carry the actual force of law, the EU AI Act's enforcement dates and the GPAI obligations, passed with little movement.

This intermittent, conflict-driven pattern is what matters most for what comes next: because attention to AI governance is driven by symbolic flashpoints rather than the legal calendar, the steps that actually bind developers reach the public with the least scrutiny, and the quiet around a compliance deadline should not be mistaken for the absence of consequence.
On the same logic, future regulatory milestones will likely surface in the news only when they turn contested or coincide with a salient capability shock.

Future work should annotate a stratified sample with LLM-based labels to estimate the classifier's precision and recall~\cite{gilardi_chatgpt_2023,ziems_can_2024}; extend the lexicon dynamically via methods such as LGDE~\cite{schindler_lgde_2025}; and track governance coverage beyond mid-2026.


\begin{acks}
During the preparation of this work, the author used \textbf{Writefull}
(Overleaf integration) in order to: \textbf{Grammar and spelling check},
\textbf{Paraphrase and reword}, \textbf{Improve writing style}.
After using these tools, the author reviewed and edited the content as
needed and takes full responsibility for the publication's content.
\end{acks}

%%
%% The next two lines define the bibliography style to be used, and
%% the bibliography file.
\bibliographystyle{ACM-Reference-Format}
\bibliography{references}

\end{document}
\endinput
%%
%% End of file `paper.tex'.
